% Fragmento integrable. Preambulo anfitrion: amsmath, amssymb, longtable, hyperref.
\section{Compatibilidad de lectores y restitución del parámetro}
\label{hmtfg:fragmento:compatibilidad-de-lectores-y-restitucion-del-parametro}

\subsection{Compatibilidad de lectores, mínimos y recuperación de la historia}
\label{hmtfg:escalado:13}

Las operaciones de prolongación y los lectores ordenados conservan la información requerida por la selección de raíces. Las siguientes identidades fijan sus dominios y su composición.

\subsubsection{Regla forward recuperada}
\label{hmtfg:escalado:13.1}

La actualización determinista parte del estado de frontera R36 con carácter \(\chi\) y residuo interno. La actualización es
\begin{equation}
\mathcal U_K^\chi=
\operatorname{Upd}_K^\chi\circ
\operatorname{Tra}_K^\chi\circ
\operatorname{Sel}_K^\chi,
\qquad
N_{K+1}=729N_K+d_K^\chi.
\label{hmtfg:escalado:eq:26}
\end{equation}
El dígito y la actualización dependen del estado precedente. Dos secciones forward con el mismo estado completo R36 coinciden por inducción. La subrelación \(\operatorname{Ext}^{\chi,\rightarrow}\) es el grafo de esa actualización; la relación total \(\operatorname{Ext}\) admite multiplicidad de continuaciones según los datos y caracteres.

Las prolongaciones compatibles transportan los lectores regionales y la firma conjunta sobre los mismos prefijos.

Estas reglas transportan orientación, residuo, cociente, acarreo, frontera y memoria. Su aplicación a la cascada debe conservar además la selección ordenada del parámetro del lector crítico.

\subsubsection{Independencia de las constantes antecedentes respecto del parámetro libre}
\label{hmtfg:escalado:13.2}

Sea \(B\) el estado HMT que proporciona la rueda y sus constantes antecedentes. Escribamos \(\mathcal C(B)\) para el conjunto de esas publicaciones y \(u_B\) para el perfil crítico. En el dominio del lector crítico, la construcción posterior es
\begin{equation}
(B,\lambda)\longmapsto f_{B,\lambda}(x)=1-\lambda u_B(x),
\qquad 0<\lambda\le 2/u_B(1).
\label{hmtfg:escalado:eq:27}
\end{equation}
Para un mismo \(B\), la lectura de las constantes anteriores se factoriza por la proyección al primer componente:
\begin{equation}
\widetilde{\mathcal C}(B,\lambda)
=\mathcal C\circ\operatorname{pr}_B(B,\lambda)
=\mathcal C(B).
\label{hmtfg:escalado:eq:28}
\end{equation}
En consecuencia, dos valores distintos de \(\lambda\) conservan las mismas publicaciones antecedentes. La igualdad es una propiedad de esta construcción causal: el perfil se fija antes de variar \(\lambda\), y el factor \(\pi_{\mathrm{HMT}}\) se cancela en su normalización cuadrática.

\textbf{Lema.} Supóngase que una condición \(A\) se evalúa únicamente sobre esas publicaciones antecedentes. Para todo \(B\) fijo,
\begin{equation}
\{\lambda\in\Lambda:A(\widetilde{\mathcal C}(B,\lambda))\}
=
\begin{cases}
\Lambda,&A(\mathcal C(B))\text{ se cumple},\\
\varnothing,&A(\mathcal C(B))\text{ falla}.
\end{cases}
\label{hmtfg:escalado:eq:29}
\end{equation}
\textbf{Demostración.} Sustituir \eqref{hmtfg:escalado:eq:28} hace independiente de \(\lambda\) el predicado. Por tanto, lo satisfacen todos los parámetros del dominio o ninguno. Esto incluye una comparación de las constantes antecedentes con intervalos metrológicos.

La unicidad de la sección \(B\) y la restricción posterior \(f_{B,\lambda}^{2^n}(0)=0\) producen una fibra de soluciones en \(\lambda\). La regla de primera raíz actúa sobre esta fibra mediante el retorno y su itinerario. La selección y su prolongación se demuestran en los resultados siguientes.

Este resultado mantiene el orden HMT: la metrología contrasta salidas previamente construidas y la selección de la raíz se demuestra en el lector que la produce.

\subsubsection{Lema de transporte ordenado del selector}
\label{hmtfg:escalado:13.3}

Sea \(P_n\) un conjunto ordenado de candidatos genealógicos con mínimo \(p_n^*\). Sea \(Z_n\) el conjunto original de raíces candidatas del corpus, con el orden real, y sea \(L_n:P_n\to Z_n\) un lector monótono. Se dice que su imagen es \textbf{coinicial} cuando
\begin{equation}
\forall z\in Z_n\quad\exists p\in P_n:
L_n(p)\le z.
\label{hmtfg:escalado:eq:30}
\end{equation}
Entonces
\begin{equation}
\boxed{L_n(p_n^*)=\min Z_n
\quad\Longleftrightarrow\quad
L_n(P_n)\text{ es coinicial en }Z_n.}
\label{hmtfg:escalado:eq:31}
\end{equation}
\textbf{Demostración.} Si la imagen del mínimo es el mínimo global, se toma \(p=p_n^*\) en \eqref{hmtfg:escalado:eq:30}. Recíprocamente, dado \(z\in Z_n\), la coinicialidad proporciona \(p\) con \(L_n(p)\le z\). La monotonía da \(L_n(p_n^*)\le L_n(p)\le z\). Como \(L_n(p_n^*)\in Z_n\), es el mínimo global. La sobreyectividad de \(L_n\) constituye una condición suficiente, más fuerte que \eqref{hmtfg:escalado:eq:30}.

El isomorfismo ordenado de la carta del continuo satisface la sobreyectividad del lector. La aplicación concreta al retorno y a su conjunto completo de ceros se demuestra en la sección de transporte ordenado; la clasificación y el aislamiento determinan después su continuación dinámica.

\subsubsection{Selección mínima mediante supervivencia a profundidad arbitraria}
\label{hmtfg:escalado:13.6}

La composición con la supervivencia produce además un selector global preciso. Se mantiene la familia uniforme \(\eta=0\), construida en sección~\ref{hmtfg:escalado:1}, y se utilizan los teoremas de realización del itinerario infinito y conservación de sus signos. Sean
\[
I=\left[\frac1{2u_0(1)},\frac2{u_0(1)}\right],\qquad
c_j(\lambda)=f_\lambda^j(0),\qquad
\tau_j=(-1)^{v_2(j)}.
\]
El conjunto
\[
\Lambda_\tau=\{\lambda\in I:\operatorname{sign}c_j(\lambda)=\tau_j
\text{ para todo }j\ge1\}
\]
es compacto y no vacío. El control de derivadas de ese propietario proporciona
\begin{equation}
B_j=\frac{16^j(16^j-1)}{15},\qquad
|c_j(\lambda)|>\frac2{B_j}\quad(\lambda\in\Lambda_\tau).
\label{hmtfg:escalado:eq:32}
\end{equation}
Definamos
\begin{equation}
\varepsilon_j=\frac1{B_j},\qquad
A_N=\{\lambda\in I:\tau_jc_j(\lambda)\ge\varepsilon_j,\
1\le j\le N\}.
\label{hmtfg:escalado:eq:33}
\end{equation}
Cada \(A_N\) es compacto por continuidad de los iterados y no vacío porque contiene \(\Lambda_\tau\). Además,
\begin{equation}
A_{N+1}\subseteq A_N,\qquad
\bigcap_{N\ge1}A_N=\Lambda_\tau.
\label{hmtfg:escalado:eq:34}
\end{equation}
La inclusión de derecha a izquierda se sigue de \eqref{hmtfg:escalado:eq:32}. La inclusión inversa se sigue de los márgenes estrictamente positivos de \eqref{hmtfg:escalado:eq:33}, que fijan cada signo.

\textbf{Teorema del mínimo superviviente.} Los mínimos \(a_N=\min A_N\) satisfacen
\begin{equation}
\boxed{a_N\uparrow a_*=\min\Lambda_\tau.}
\label{hmtfg:escalado:eq:35}
\end{equation}
\textbf{Demostración.} La inclusión \eqref{hmtfg:escalado:eq:34} hace creciente la sucesión de mínimos, acotada dentro de \(I\); tiene, por tanto, un límite \(a\). Para cada \(r\), todos los términos \(a_N\), \(N\ge r\), pertenecen al cerrado \(A_r\); así \(a\in A_r\). Por \eqref{hmtfg:escalado:eq:34}, \(a\in\Lambda_\tau\). Para cualquier \(\lambda\in\Lambda_\tau\), se tiene \(a_N\le\lambda\) a todo nivel, y al pasar al límite, \(a\le\lambda\). Esto prueba \eqref{hmtfg:escalado:eq:35}.

En el lector cilíndrico HMT, el objeto \(a_*\) posee prefijos compatibles a toda profundidad, con orientación, residuo y memoria conservados. Su elección utiliza el orden de la lectura paramétrica y la supervivencia completa. La compacidad demuestra la existencia y convergencia de estos mínimos.

Los parámetros \(a_N\) resuelven restricciones de signos con margen. Sus fronteras pueden satisfacer \(\tau_jc_j=\varepsilon_j\), mientras las raíces superestables satisfacen \(c_{2^n}=0\). Por ello \eqref{hmtfg:escalado:eq:35} establece un selector global de realización infinita y preserva, como objeto distinto, el selector de primeras raíces del corpus.

\subsubsection{Recuperación del parámetro desde la historia crítica completa}
\label{hmtfg:escalado:13.7}

El segundo término de la historia crítica proporciona un lector inverso exacto:
\begin{equation}
c_1(\lambda)=1,\qquad
c_2(\lambda)=1-\lambda u_0(1),\qquad
\boxed{\lambda=\frac{1-c_2(\lambda)}{u_0(1)}.}
\label{hmtfg:escalado:eq:36}
\end{equation}
La positividad \(u_0(1)>0\), demostrada para el perfil, hace válida la división. En consecuencia, dos historias críticas completas iguales tienen el mismo parámetro. La palabra de signos \(\tau\) constituye una publicación reducida de esa historia; conserva su combinatoria y omite la magnitud de \(c_2\).

Las fórmulas \eqref{hmtfg:escalado:eq:35}–\eqref{hmtfg:escalado:eq:36} unen la selección por supervivencia y la restitución del parámetro desde la historia crítica. La igualdad de su límite con el cruce estable se demuestra por primera salida, y la identificación de los centros por el transporte analítico de sus clases híbridas.
